\documentclass{../../../unipd-notes}

\unipdsetup{
  course = {Fondamenti di Ingegneria Informatica},
  short-course = {Fondamenti di Ingegneria},
  document-type = {Appunti delle lezioni}
}

\begin{document}
\chapter{Figure e tabelle}

\begin{figure}[htbp]
  \centering
  \fcolorbox{unipd-rule}{unipd-soft}{%
    \parbox[c][32mm][c]{0.68\textwidth}{\centering\sffamily
      Fotografia della scheda di sviluppo\\[2mm]
      \footnotesize STM32 Nucleo utilizzata in laboratorio}%
  }
  \caption{Scheda di sviluppo impiegata nelle esercitazioni}
  \label{fig:control-loop}
  \figuresource{elaborazione propria}
\end{figure}

La \cref{fig:control-loop} identifica l'hardware di riferimento per le prove di acquisizione analogica.

\begin{table}[htbp]
  \centering
  \caption{Parametri elettrici misurati sul convertitore analogico-digitale}
  \label{tab:measurements}
  \begin{tabularx}{0.72\textwidth}{LCC}
    \toprule
    Grandezza & Valore & Unità \\
    \midrule
    Tensione di riferimento & $3.30$ & \si{\volt} \\
    Corrente assorbita & $18.4$ & \si{\milli\ampere} \\
    Frequenza di campionamento & $10.0$ & \si{\kilo\hertz} \\
    \bottomrule
  \end{tabularx}
  \tablesource{misure raccolte durante l'esercitazione di laboratorio}
\end{table}

I valori della \cref{tab:measurements} sono stati acquisiti con alimentazione USB e carico nominale.
\end{document}
